% =========================================================
% Community-Aware Early Warning of Distrust Emergence
% in Temporal Bitcoin Trust Networks
%
% IEEE conference paper draft.
% Compile with pdflatex; uses IEEEtran.cls (IEEEtran v1.8b or later).
% =========================================================
\documentclass[conference]{IEEEtran}

\IEEEoverridecommandlockouts
% packages
\usepackage{cite}
\usepackage{amsmath,amssymb,amsfonts}
\usepackage{algorithmic}
\usepackage{graphicx}
\usepackage{textcomp}
\usepackage{xcolor}
\usepackage{booktabs}
\usepackage{url}
\usepackage{hyperref}

% Avoid IEEEtran title-style font conflicts
\def\BibTeX{{\rm B\kern-.05em{\sc i\kern-.025em b}\kern-.08em
    T\kern-.1667em\lower.7ex\hbox{E}\kern-.125emX}}

% Inline code shortcut
\newcommand{\code}[1]{\texttt{#1}}

% TODO marker
\newcommand{\todo}[1]{\textcolor{red}{\textbf{[TODO: #1]}}}

\begin{document}

% =========================================================
% TITLE BLOCK
% =========================================================
\title{Community-Aware Early Warning of Distrust Emergence \\
in Temporal Bitcoin Trust Networks}

\author{
\IEEEauthorblockN{Author Name}
\IEEEauthorblockA{\textit{Department / Institution} \\
\textit{University / Affiliation}\\
City, Country \\
email@institution.edu}
}
% \todo{fill in author block before submission}

\maketitle

% =========================================================
% ABSTRACT
% =========================================================
\begin{abstract}
Bitcoin-OTC was a peer-rated trust marketplace whose users rated each
other on a $[-10, +10]$ scale, producing a signed, weighted, directed,
timestamped social network. We study early warning of \emph{new
negative} trust edges in this network: at the end of month $t$, can we
rank ordered pairs of users by their risk of receiving a new negative
directed rating in month $t+1$, using only information available at
time $t$? We build an interpretable, leakage-controlled pipeline with
cumulative monthly snapshots and a locked train\,/\,validation\,/\,test
split, and run a three-way feature ablation
(structural; structural + temporal;
structural + temporal + community-aware) under Logistic Regression and
Random Forest. Communities are extracted with Louvain on the
positive-edge projection of each snapshot. On a test fold of 34
positives over 148{,}639 candidate pairs, structural features alone
already produce a non-trivial ranker (Random Forest ROC-AUC $0.88$,
PR-AUC $0.022$). Adding temporal features delivers the largest gain;
the $(\mathrm{S}{+}\mathrm{T}) - \mathrm{S}$ ROC-AUC delta is robust
under a paired stratified bootstrap (95\% CI well above zero, 100\% of
resamples positive). Adding community-aware features on top of
temporal features does not improve ROC-AUC, but shifts Random Forest
PR-AUC upward by a directionally positive but not statistically
decisive amount (mean $+0.026$, 94\% of resamples positive, 95\% CI
just including zero). We therefore report community gains as
suggestive rather than conclusive at this test-set size, and identify
\emph{pair-level neighborhood overlap} -- rather than a coarse
same-community vs cross-community indicator -- as the most useful
community-aware signal.
\end{abstract}

\begin{IEEEkeywords}
signed networks, temporal link prediction, community detection,
Louvain, Random Forest, bootstrap, Bitcoin-OTC, weak ties.
\end{IEEEkeywords}

% =========================================================
\section{Introduction}
% =========================================================

Online trust marketplaces let participants rate one another so that
newcomers can assess the reliability of potential counterparties.
Bitcoin-OTC, an over-the-counter cryptocurrency marketplace, hosted
such a web-of-trust between 2010 and early 2016. Each rating is a
directed signed weighted timestamped event: user $u$ assigns user $v$
an integer rating in $[-10, +10]$, and the rating is stored with its
timestamp. The resulting graph is one of the canonical benchmarks for
signed network analysis~\cite{kumar2016}.

This paper studies \emph{early warning of distrust emergence}: given
the network at the end of month $t$, we want to rank ordered pairs
$(u, v)$ by the probability that $v$ will receive a new \emph{negative}
directed rating from $u$ in month $t+1$. Useful applications include
marketplace moderation, fraud surveillance, and reputation governance.
The task differs from generic link prediction in three ways. First,
negative ratings are rare: only about 10\% of edges in Bitcoin-OTC are
negative, and new negative events per month are far rarer still.
Second, the task is sign-dependent: predicting any new edge and
predicting a new negative edge are different problems, and positive-edge
popularity heuristics do not directly transfer. Third, the evaluation
must respect the arrow of time -- a model that uses post-$t$
information would not be deployable for warning.

Our research question is:
\begin{quote}
At snapshot $t$, can we predict whether an ordered pair $(u, v)$ --
either previously unconnected, or connected only positively -- will
receive a new negative directed edge in the window
$(t, t{+}1]$, using only information available at time $t$?
\end{quote}

Our contributions are:
\begin{itemize}
\item An interpretable feature-based pipeline -- classical
link-prediction, signed structural, trust-summary, temporal, and
community-aware features -- with explicit leakage controls and a
locked time-based train\,/\,validation\,/\,test split.
\item A three-way feature ablation (structural; structural + temporal;
structural + temporal + community) under Logistic Regression and
Random Forest, evaluated on a locked test fold.
\item Paired stratified bootstrap confidence intervals for all
reported test metrics and for the metric \emph{differences} between
feature sets, so that observed improvements are reported with sampling
uncertainty rather than as point estimates.
\item A descriptive analysis of community structure and weak-tie
signals in the context of distrust emergence.
\end{itemize}

The pipeline is intentionally lighter than the deep signed-graph
methods surveyed by Zhang \emph{et al.}~\cite{zhang2024}. Our goal is
interpretable analysis on a single dataset; we do not claim
state-of-the-art absolute performance.

% =========================================================
\section{Related Work}
% =========================================================

\textbf{Signed networks and weighted trust prediction.} The
Bitcoin-OTC and Bitcoin-Alpha datasets, together with the iterative
\emph{fairness} and \emph{goodness} statistics for nodes, were
introduced by Kumar \emph{et al.}~\cite{kumar2016}. We do not
reimplement iterated fairness/goodness here; instead, we use four
lightweight per-node trust-summary features (mean rating given, mean
rating received, negative-given ratio, negative-received ratio) as a
deliberately simple analogue.

\textbf{Behavioral dynamics on Bitcoin-OTC.} Bertazzi
\emph{et al.}~\cite{bertazzi2018} analyze the multilayer structure of
the Bitcoin-OTC web of trust, contrasting rewarding and punitive
behaviors at the network level. This motivates our use of cumulative
monthly snapshots and our temporal\,/\,recency features (days since
last activity, recent-window counts, and exponentially decayed
aggregates).

\textbf{Community-aware temporal prediction.}
Choudhury~\cite{choudhury2024} develops community-aware dynamic
similarity features for supervised link prediction on dynamic
networks, arguing that community structure carries predictive signal
beyond pair-level topology. We operationalize this idea by running
Louvain on the positive-edge projection of every snapshot and
engineering community-aware pair features
(\code{same\_community}, community sizes, \code{neighborhood\_overlap},
and boundary fractions).

\textbf{Deep signed graph learning.} Zhang \emph{et al.}~\cite{zhang2024}
survey signed graph representation learning, including
GNN-based approaches such as signed graph convolutional and attention
networks. We use this body of work for positioning only: our approach
is shallow and interpretable, and we do not benchmark against GNNs.

\textbf{Classical heuristics and structural theory.} We use Common
Neighbors, Jaccard similarity, Adamic-Adar~\cite{adamic2003}, and
Preferential Attachment on the undirected any-sign projection of the
snapshot graph. Our community-aware features include the
neighborhood-overlap diagnostic associated with Granovetter's
strength-of-weak-ties hypothesis~\cite{granovetter1973}. Louvain
community detection follows Blondel \emph{et al.}~\cite{blondel2008}.

% =========================================================
\section{Dataset and Problem Formulation}
% =========================================================

\subsection{Bitcoin-OTC and basic terminology}
Each row of the public dataset (file \code{soc-sign-bitcoinotc.csv},
SNAP~\cite{snap}) is a single rating event with fields
\code{(source, target, rating, timestamp)}. A \emph{node} is a user.
A \emph{directed edge} $u \to v$ is a rating event from $u$ to $v$. The
\emph{sign} of an edge is $+1$ if $\text{rating} > 0$ and $-1$ if
$\text{rating} < 0$. The \emph{weight} of an edge is the integer
rating in $[-10, +10]$. The \emph{timestamp} is a Unix-seconds float;
the dataset spans 2010-11-08 to 2016-01-25, approximately 63 calendar
months.

After defensive cleaning (drop self-loops and zero ratings; neither
occurs in this file), the corpus has 35{,}592 rating events, 5{,}881
distinct users, 32{,}029 positive and 3{,}563 negative ratings
($\approx 9{:}1$ imbalance), and reciprocity $0.79$. The giant
weakly connected component has 5{,}875 nodes and the giant strongly
connected component has 4{,}709 nodes. Fig.~\ref{fig:edges_per_month}
shows the monthly interaction volume; Fig.~\ref{fig:rating_histogram}
shows the rating distribution.

\begin{figure}[t]
\centering
\includegraphics[width=\columnwidth]{../results/figures/edges_per_month.png}
\caption{Monthly interaction volume on Bitcoin-OTC, Nov-2010 through
Jan-2016. Activity grows through 2014--2015 and then collapses; the
tail collapse motivates the eligibility filter discussed in
Section~\ref{sec:setup}.}
\label{fig:edges_per_month}
\end{figure}

\begin{figure}[t]
\centering
\includegraphics[width=\columnwidth]{../results/figures/rating_histogram.png}
\caption{Distribution of rating values across all events. Positive
ratings (especially $+1$ and $+10$) dominate; negatives are a long left
tail.}
\label{fig:rating_histogram}
\end{figure}

\subsection{Cumulative monthly snapshots}
We adopt a cumulative monthly schedule. For each month index $t \in
\{0, \dots, 62\}$, the snapshot graph at time $t$ contains all events
with timestamp at or before the end of month $t$. When the pair
$(u, v)$ appears multiple times in the cumulative window, we collapse
the multi-edge to its \emph{latest} rating before the end of month
$t$. This encodes ``$u$'s current trust judgment of $v$'' and is used
\emph{only for the snapshot graph state}. Temporal and history-based
features instead use the full event history with
$\mathrm{ts} \leq \mathrm{end\_of\_month}(t)$, preserving the
underlying multi-event structure.

From the directed snapshot graph $G^t_{\mathrm{dir}}$ we derive two
undirected weighted projections: $G^t_{\mathrm{pos}}$ keeps only edges
whose collapsed sign is $+1$, with weight equal to the number of
distinct directed positive contributions; $G^t_{\mathrm{neg}}$ is the
analogous negative projection.

\subsection{The task}
Given snapshot $t$, the target is a binary label per ordered candidate
pair $(u, v)$:
\begin{equation*}
y(u, v) =
\begin{cases}
1, & \exists\,(u, v, r, \mathrm{ts}) \text{ with } r < 0 \\
   & \text{and } \mathrm{end\_of\_month}(t) < \mathrm{ts} \\
   & \quad\leq \mathrm{end\_of\_month}(t{+}1), \\
0, & \text{otherwise.}
\end{cases}
\end{equation*}
Both ``no new $u \to v$ event'' and ``only new positive $u \to v$
events'' map to $y = 0$; only a new \emph{negative} directed rating
counts. This includes the case where $u$ previously rated $v$
positively but flips to negative in the next window, which is a
legitimate form of distrust emergence and is therefore retained as a
positive example.

\subsection{Why ranking metrics rather than thresholded
precision/recall}
The candidate set for any non-trivial snapshot is large (approximately
$1 \times 10^5$ to $2 \times 10^5$ ordered pairs after filtering;
see Section~\ref{sec:method}), and the prior of $y = 1$ is on the
order of $10^{-4}$. At this imbalance, a threshold-$0.5$ classifier is
not a meaningful operating point. A model with
\code{class\_weight = 'balanced'} tends to predict the positive class
at unrealistic precision; a default Random Forest predicts essentially
no positives. Threshold-dependent precision, recall, and F1 are
therefore nearly uninformative in this regime. We focus on
\emph{ranking} metrics: ROC-AUC, PR-AUC (average precision), and
Precision@$k$ for small $k$. These map naturally onto an operational
watchlist of the $k$ riskiest pairs.

% =========================================================
\section{Methodology}\label{sec:method}
% =========================================================

\subsection{Core node filtering}
A node $n$ is \emph{core at snapshot $t$} iff (a) its total
interactions (in + out events) up to $\mathrm{end\_of\_month}(t)$ is at
least \code{MIN\_ACTIVITY = 3}, and (b) its most recent activity is
within \code{RECENT\_DAYS = 180} of $\mathrm{end\_of\_month}(t)$. Both
thresholds are configuration parameters; we adopt these values as
starting defaults and do not sweep them.

\subsection{Candidate pair generation}
For each snapshot $t$, candidate ordered pairs $(u, v)$ satisfy:
(1)~$u, v \in \mathrm{Core}(t)$ and $u \neq v$;
(2)~$v$ is within \code{K\_HOP = 2} hops of $u$ on the undirected
any-sign projection of $G^t_{\mathrm{dir}}$;
(3)~there is no existing \emph{negative} directed edge $u \to v$ at
time $t$ -- pairs with a prior \emph{positive} $u \to v$ are retained,
because a rating flip from positive to negative in the next window is
itself distrust emergence;
(4)~the total candidate count is capped at
\code{MAX\_CANDIDATES = 200{,}000} via fixed-seed uniform
downsampling. The cap binds only at the densest mid-period snapshots.

\subsection{Label generation and leakage control}
Labels are derived from the raw event history in the window
$(\mathrm{end\_of\_month}(t), \mathrm{end\_of\_month}(t{+}1)]$. A
defensive assertion in every feature precompute confirms that the
maximum timestamp used during feature computation does not exceed
$\mathrm{end\_of\_month}(t)$.

\subsection{Feature families}
The full feature vector has 36 columns split across five families:

\textbf{Classical link prediction (4 features).} Common Neighbors,
Jaccard similarity, Adamic-Adar~\cite{adamic2003}, Preferential
Attachment, all computed on the undirected any-sign projection of
$G^t_{\mathrm{dir}}$.

\textbf{Signed structural (8 features).} Positive and negative
out-degree of $u$; positive and negative in-degree of $v$; positive
and negative common-neighbor counts (on $G^t_{\mathrm{pos}}$ and
$G^t_{\mathrm{neg}}$ respectively); reciprocity flag (1 if $v \to u$
exists in $G^t_{\mathrm{dir}}$); prior-interaction flag (1 if
$u \to v$ exists in $G^t_{\mathrm{dir}}$ -- necessarily a positive
prior, by construction of the candidate set, so this flag is partly
degenerate and we report it transparently).

\textbf{Trust summary (4 features).} Mean rating given by $u$; mean
rating received by $v$; fraction of $u$'s given ratings that were
negative; fraction of $v$'s received ratings that were negative.
These are a deliberately simple, non-iterated analogue of the
fairness/goodness statistics in~\cite{kumar2016}.

\textbf{Temporal / recency (12 features).} Days since $u$'s last
activity; days since $v$'s last activity; days since the last
$u \to v$ interaction (sentinel value $9999$ when no prior $u \to v$
event exists); count of past $u \to v$ events; recent-window (90-day)
counts of negatives given by $u$, negatives received by $v$, total
events given by $u$, total events received by $v$; and four
exponentially decayed aggregates with half-life $\tau = 60$ days
(decayed total activity and decayed negative count, for both $u$ and
$v$).

\textbf{Community-aware (8 features).} \code{same\_community}
indicator; its complement \code{cross\_community\_pair};
\code{community\_size\_u}, \code{community\_size\_v}; the log-size
gap $|\log s_u - \log s_v|$;
\code{neighborhood\_overlap} $= |N(u) \cap N(v) \setminus \{u,v\}|
/ |N(u) \cup N(v) \setminus \{u,v\}|$ on the undirected any-sign
projection; and \code{boundary\_frac\_u}, \code{boundary\_frac\_v},
each defined for a node $n$ as the fraction of $n$'s neighbors in
$G^t_{\mathrm{pos}}$ whose Louvain community differs from $n$'s.

\subsection{Community extraction}\label{sec:method:community}
For each snapshot $t$ we run Louvain~\cite{blondel2008} on the
positive-edge undirected weighted graph $G^t_{\mathrm{pos}}$. Louvain
is stochastic; we fix \code{random\_state} and cache the resulting
partition per snapshot so that downstream community-feature
extraction is deterministic. Across the 37 split snapshots (see
Section~\ref{sec:setup}) modularity sits in $[0.47, 0.52]$ with a slight
upward trend, while the number of communities grows from $\sim 17$ at
$t = 13$ to $\sim 40$--$47$ at $t = 49$ (Fig.~\ref{fig:modularity}).

\begin{figure}[t]
\centering
\includegraphics[width=\columnwidth]{../results/figures/modularity_over_time.png}
\caption{Louvain modularity (left axis) and number of communities
(right axis) on the per-snapshot positive-edge graph
$G^t_{\mathrm{pos}}$. Modularity is stable in $[0.47, 0.52]$; the
community count grows roughly linearly as the network grows. Dotted
lines mark the train\,/\,val and val\,/\,test boundaries.}
\label{fig:modularity}
\end{figure}

\subsection{Models}
We train three models per feature set: (i) a \emph{random baseline}
that draws uniform probabilities (seeded), to calibrate the floor on
each metric; (ii) \emph{Logistic Regression} with \code{StandardScaler}
and \code{class\_weight = 'balanced'}, \code{liblinear} solver,
\code{max\_iter = 2000}; (iii) \emph{Random Forest} with
\code{n\_estimators = 200}, \code{min\_samples\_leaf = 20},
\code{class\_weight = 'balanced\_subsample'}. We perform no
hyperparameter tuning beyond these defaults; the focus is feature
engineering and ablation.

\subsection{Training-set negative subsampling}
The training split contains $\sim 5 \times 10^{6}$ candidate pairs and
only 545 positives. To keep training tractable while preserving all
positives, we subsample negatives in TRAIN to a $100{:}1$
negative:positive ratio, yielding $\sim 55{,}000$ training rows.
Validation and test sets remain complete; this is a training
acceleration, not an evaluation trick. A fixed RNG seed makes the
subsample identical across feature sets and models, so paired
comparisons within a model are valid.

\subsection{Evaluation metrics}
We report ROC-AUC, PR-AUC (average precision), and Precision@$k$ for
$k \in \{50, 100\}$. We also record thresholded precision, recall, and
F1 at the default $0.5$ cutoff but read them as uninformative at this
imbalance.

% =========================================================
\section{Experimental Setup}\label{sec:setup}
% =========================================================

\subsection{Train / validation / test split}
A snapshot is \emph{eligible} to enter the split if (i) its next
window contains at least \code{MIN\_POS\_FOR\_SPLIT = 10} raw new
negative events, and (ii) it sits in the longest contiguous block of
eligible snapshots (this excludes a sporadic late-dataset island at
$t = 53$). The resulting eligible pool is $t \in \{13, \dots, 49\}$
(37 snapshots). Within the pool we apply a forward-time split: the
last two snapshots are TEST, the previous two are VALIDATION, and the
remainder is TRAIN. The split is locked from the start of feature
engineering and is never re-tuned. Exact sizes are in
Table~\ref{tab:split}.

\begin{table}[t]
\caption{Locked train / validation / test split sizes.}
\label{tab:split}
\centering
\begin{tabular}{lccc}
\toprule
Split & Snapshots & Pairs & Positives \\
\midrule
TRAIN & $t = 13 \dots 45$ & 5{,}013{,}346 & 545 \\
VAL   & $t = 46, 47$      &   184{,}648   &  16 \\
TEST  & $t = 48, 49$      &   148{,}639   &  34 \\
\bottomrule
\end{tabular}
\end{table}

The small positive counts in VAL and TEST are an unavoidable
consequence of the dead tail of Bitcoin-OTC in late 2015. We discuss
this explicitly in Section~\ref{sec:lim}.

\subsection{Bootstrap protocol}
We attach uncertainty estimates to all final test metrics with a
\emph{paired stratified bootstrap}. For $B = 1000$ iterations, we draw
34 positive indices with replacement and 148{,}605 negative indices
with replacement (preserving the original positive prevalence and
guaranteeing both labels in every resample). The same resample index
vector is reused across all three feature sets within one model, so
bootstrap deltas such as
$(\mathrm{S}{+}\mathrm{T}{+}\mathrm{C}) - (\mathrm{S}{+}\mathrm{T})$
are properly paired. Random Forest produces probability scores on a
discrete $1/200$ grid, so resampling with replacement creates ties at
the top of the score list. To prevent \code{argpartition} tie order
from biasing Precision@$k$, we add a $\mathcal{U}[0, 10^{-12}]$ jitter
to probabilities before ranking, which leaves non-tied rankings
unchanged and randomizes ties. We report 95\% percentile confidence
intervals and the bootstrap-positive share $\mathrm{frac\_boot}(>0)$;
this share is the proportion of bootstrap iterations in which the
delta is strictly positive and is \emph{not} a frequentist p-value.
We interpret a CI that overlaps zero as ``not distinguishable from
sampling noise on this test set.''

\subsection{Reproducibility and artifacts}
All seeds (model fits, training subsample, Louvain partition,
bootstrap) are fixed at $42$. Intermediate parquet files and a
canonical test-set probability matrix are persisted so that
downstream analysis can be re-run without re-training. Code is
organized in \code{src/} (algorithmic modules) and \code{scripts/}
(stage entry points). The headline numbers in the tables below come
from \code{results/test\_metrics\_from\_predictions.csv}, which
recomputes every metric directly from the saved test parquet so that
the headline table, figure legends, and bootstrap analyses all derive
from the same probability vectors.\footnote{An older
\code{stage7\_test\_metrics.csv} is also on disk and is kept as a
historical record. It differs from the canonical table by at most
$\sim 0.01$ PR-AUC for the structural-only Random Forest row; the
drift is consistent with float-level non-determinism in
\code{RandomForestClassifier.predict\_proba} with \code{n\_jobs = -1}
between two independent fits and is well inside the bootstrap CI for
that cell.}

% =========================================================
\section{Results}\label{sec:results}
% =========================================================

We label feature sets S (structural + trust summary, 16 columns),
S$+$T (+temporal, 28 columns), and S$+$T$+$C (+community, 36 columns).
Table~\ref{tab:test_metrics} contains the canonical test-set headline.

\begin{table}[t]
\caption{Final test-set metrics (test = $t \in \{48, 49\}$;
34 positives, 148{,}639 candidate pairs).}
\label{tab:test_metrics}
\centering
\begin{tabular}{llcccc}
\toprule
F.\,set & Model & ROC-AUC & PR-AUC & P@50 & P@100 \\
\midrule
S      & logreg & 0.909 & 0.0034 & 0.00 & 0.01 \\
S$+$T   & logreg & 0.949 & 0.0059 & 0.00 & 0.01 \\
S$+$T$+$C & logreg & 0.945 & 0.0058 & 0.00 & 0.01 \\
\addlinespace
S      & rf     & 0.880 & 0.0223 & 0.04 & 0.03 \\
S$+$T   & rf     & 0.950 & 0.0400 & 0.08 & 0.07 \\
S$+$T$+$C & rf     & 0.940 & 0.0703 & 0.10 & 0.06 \\
\bottomrule
\end{tabular}
\end{table}

% ----------------------------------------------------------
\subsection{Structural baseline}
% ----------------------------------------------------------
Even with structural and trust-summary features alone, both models
clearly beat random ranking. Random Forest reaches ROC-AUC $0.880$ and
PR-AUC $0.0223$, almost two orders of magnitude above the base rate of
$2.3 \times 10^{-4}$. Logistic Regression reaches ROC-AUC $0.909$ and
PR-AUC $0.0034$. The two models trade off differently: LR has higher
ROC-AUC at this feature set, while RF has higher PR-AUC and stronger
top-of-list precision. We read this as evidence that classical
heuristics and simple signed-degree summaries already carry
meaningful early-warning signal within this dataset.

% ----------------------------------------------------------
\subsection{Temporal ablation}
% ----------------------------------------------------------
Adding the 12 temporal features produces the largest single jump in
the ablation. Random Forest ROC-AUC moves from $0.880$ to $0.950$ and
PR-AUC roughly doubles from $0.0223$ to $0.0400$; Precision@50 doubles
($0.04 \to 0.08$). Logistic Regression also improves on ranking
($0.909 \to 0.949$ ROC-AUC) but its absolute PR-AUC remains small.

Section~\ref{sec:results:bootstrap} shows that the ROC-AUC delta from
temporal features is \emph{robust} under bootstrap, while the PR-AUC
and Precision@$k$ deltas are positive in direction but their 95\%
CIs touch zero. We therefore read temporal features as the main
predictive driver in this pipeline, with the strongest evidence on
ROC-AUC.

The temporal feature family is internally redundant by design --
recent-window counts and exponentially decayed aggregates carry
overlapping information, with within-pair correlations $\geq 0.9$ for
some pairs. We retain both to keep the feature set transparent rather
than to maximize parsimony.

% ----------------------------------------------------------
\subsection{Community-aware ablation}
% ----------------------------------------------------------
On top of structural and temporal features, the community-aware family
shows a metric-specific pattern (Table~\ref{tab:test_metrics}):
ROC-AUC does \emph{not} improve, while Random Forest PR-AUC increases
from $0.040$ to $0.070$ and Precision@50 from $0.08$ to $0.10$.
Logistic Regression shows essentially no movement in any metric when
community features are added.

At the descriptive level, the data are not consistent with the naive
``distrust at community boundaries'' reading. Restricting to TRAIN,
same-community pairs have an empirical positive rate of $0.0140\%$
($215 / 1{,}540{,}969$), while cross-community pairs are at $0.0095\%$
($330 / 3{,}472{,}377$): same-community pairs are \emph{more} likely
to receive a new negative rating in the next window
(Fig.~\ref{fig:samecross}).

\begin{figure}[t]
\centering
\includegraphics[width=0.85\columnwidth]{../results/figures/same_vs_cross_positive_rate.png}
\caption{TRAIN positive rate by community membership. Same-community
pairs are $\sim 47\%$ more likely to receive a new negative rating
than cross-community pairs. The naive ``cross-community distrust''
hypothesis is not supported.}
\label{fig:samecross}
\end{figure}

In the fitted Logistic Regression on S$+$T$+$C,
\code{neighborhood\_overlap} is the strongest single community feature
(standardized coefficient $-1.36$, in the direction that lower
shared-neighborhood predicts higher new-negative probability);
it is also among the top-ranked community features in Random Forest
(Table~\ref{tab:imp_lr_topfam} and
Table~\ref{tab:imp_rf_topfam}). The coarse
same/cross indicator carries little weight once
\code{neighborhood\_overlap}, boundary fractions, and community sizes
are present.

\begin{table}[t]
\caption{Top-5 features per family, Logistic Regression on S$+$T$+$C
(standardized coefficients; $+$ drives the positive class).}
\label{tab:imp_lr_topfam}
\centering
\small
\begin{tabular}{llr}
\toprule
Family & Feature & Coef. \\
\midrule
struct.\ & cn                          & $+1.28$ \\
struct.\ & prior\_interaction          & $-1.17$ \\
struct.\ & pos\_cn                     & $-1.09$ \\
struct.\ & reciprocity\_flag           & $+0.88$ \\
struct.\ & jaccard                     & $+0.88$ \\
\addlinespace
temp.\   & days\_since\_last\_activity\_u  & $-1.29$ \\
temp.\   & days\_since\_last\_uv\_interaction & $+1.19$ \\
temp.\   & uv\_interaction\_count\_past & $-1.17$ \\
temp.\   & decayed\_total\_activity\_v & $+0.76$ \\
temp.\   & decayed\_total\_activity\_u & $+0.71$ \\
\addlinespace
comm.\   & neighborhood\_overlap       & $-1.36$ \\
comm.\   & boundary\_frac\_u           & $+0.17$ \\
comm.\   & community\_size\_diff\_log  & $-0.10$ \\
comm.\   & community\_size\_v          & $-0.09$ \\
comm.\   & same\_community             & $+0.08$ \\
\bottomrule
\end{tabular}
\end{table}

\begin{table}[t]
\caption{Top-5 features per family, Random Forest on S$+$T$+$C
(Gini importance).}
\label{tab:imp_rf_topfam}
\centering
\small
\begin{tabular}{llr}
\toprule
Family & Feature & Importance \\
\midrule
struct.\ & aa                          & $0.050$ \\
struct.\ & pa                          & $0.037$ \\
struct.\ & neg\_in\_v                  & $0.034$ \\
struct.\ & reciprocity\_flag           & $0.032$ \\
struct.\ & neg\_received\_ratio\_v     & $0.030$ \\
\addlinespace
temp.\   & decayed\_total\_activity\_u & $0.109$ \\
temp.\   & days\_since\_last\_activity\_u & $0.095$ \\
temp.\   & recent\_total\_given\_u     & $0.053$ \\
temp.\   & decayed\_neg\_received\_v   & $0.051$ \\
temp.\   & decayed\_neg\_given\_u      & $0.051$ \\
\addlinespace
comm.\   & boundary\_frac\_v           & $0.018$ \\
comm.\   & neighborhood\_overlap       & $0.018$ \\
comm.\   & boundary\_frac\_u           & $0.017$ \\
comm.\   & community\_size\_v          & $0.017$ \\
comm.\   & community\_size\_u          & $0.016$ \\
\bottomrule
\end{tabular}
\end{table}

% ----------------------------------------------------------
\subsection{Test ROC and PR curves}\label{sec:results:curves}
% ----------------------------------------------------------

Fig.~\ref{fig:rf_roc} overlays Random Forest ROC curves for the three
feature sets on the test split. The S$+$T (blue) and S$+$T$+$C (red)
curves nearly coincide in the operationally relevant low-FPR region;
both are clearly above the structural-only S (grey) curve. This is the
visual analogue of the ROC saturation observed in
Table~\ref{tab:test_metrics}.

Fig.~\ref{fig:rf_pr} shows the corresponding precision-recall curves on
a logarithmic precision axis. Here the S$+$T$+$C curve sits visibly
above S$+$T between recall $\sim 0.05$ and recall $\sim 0.30$ -- the
range that matters for an early-warning watchlist. The
structural-only curve trails throughout. Average precisions (the
PR-AUC values reported in Table~\ref{tab:test_metrics}) are also
recorded in the figure legend.

\begin{figure}[t]
\centering
\includegraphics[width=\columnwidth]{../results/figures/test_rf_roc.png}
\caption{Test ROC curves for Random Forest. S$+$T and S$+$T$+$C are
visually indistinguishable in the operationally relevant low-FPR
region; both clearly dominate the structural-only S curve.}
\label{fig:rf_roc}
\end{figure}

\begin{figure}[t]
\centering
\includegraphics[width=\columnwidth]{../results/figures/test_rf_pr.png}
\caption{Test precision-recall curves for Random Forest, log
precision axis. The community-augmented S$+$T$+$C curve sits visibly
above S$+$T between recall $\sim 0.05$ and recall $\sim 0.30$,
illustrating the PR-oriented community lift; the structural-only S
curve trails throughout. Dashed line marks the base rate of
$2.29 \times 10^{-4}$.}
\label{fig:rf_pr}
\end{figure}

Logistic Regression overlays (not shown; available on disk as
\code{test\_logreg\_\{roc,pr\}.png}) tell a different story: the
S$+$T and S$+$T$+$C curves are essentially indistinguishable, while
S is below. This visually confirms the absence of any LR community
gain.

% ----------------------------------------------------------
\subsection{Bootstrap uncertainty analysis}\label{sec:results:bootstrap}
% ----------------------------------------------------------

Tables~\ref{tab:boot_abs} and~\ref{tab:boot_delta} summarize the
paired stratified bootstrap ($B = 1000$).

\begin{table}[t]
\caption{Bootstrap 95\% percentile CIs on absolute test metrics
(Random Forest).}
\label{tab:boot_abs}
\centering
\small
\begin{tabular}{lcc}
\toprule
F.\,set & ROC-AUC & PR-AUC \\
\midrule
S        & $0.88$ $[0.80, 0.94]$ & $0.030$ $[0.003, 0.097]$ \\
S$+$T    & $0.95$ $[0.92, 0.97]$ & $0.052$ $[0.011, 0.131]$ \\
S$+$T$+$C & $0.94$ $[0.90, 0.97]$ & $0.079$ $[0.012, 0.180]$ \\
\midrule
F.\,set & P@50 & P@100 \\
\midrule
S        & $0.05$ $[0.00, 0.12]$ & $0.03$ $[0.00, 0.07]$ \\
S$+$T    & $0.08$ $[0.02, 0.16]$ & $0.07$ $[0.02, 0.11]$ \\
S$+$T$+$C & $0.09$ $[0.02, 0.18]$ & $0.06$ $[0.03, 0.11]$ \\
\bottomrule
\end{tabular}
\end{table}

\begin{table}[t]
\caption{Bootstrap 95\% percentile CIs on paired metric deltas
(Random Forest). \emph{frac\_boot} is the share of resamples with
$\Delta > 0$; it is not a p-value.}
\label{tab:boot_delta}
\centering
\small
\begin{tabular}{lllc}
\toprule
$\Delta$ & Metric & 95\% CI & frac\_boot \\
\midrule
$(\mathrm{S}{+}\mathrm{T}) - \mathrm{S}$
  & ROC-AUC & $[+0.028, +0.123]$ & $1.000$ \\
  & PR-AUC  & $[-0.021, +0.082]$ & $0.849$ \\
  & P@50    & $[-0.040, +0.100]$ & $0.755$ \\
  & P@100   & $[-0.000, +0.080]$ & $0.915$ \\
\addlinespace
$(\mathrm{S}{+}\mathrm{T}{+}\mathrm{C}) - (\mathrm{S}{+}\mathrm{T})$
  & ROC-AUC & $[-0.027, +0.002]$ & $0.073$ \\
  & PR-AUC  & $[-0.003, +0.068]$ & $0.935$ \\
  & P@50    & $[+0.000, +0.060]$ & $0.465$ \\
  & P@100   & $[-0.030, +0.020]$ & $0.157$ \\
\addlinespace
$(\mathrm{S}{+}\mathrm{T}{+}\mathrm{C}) - \mathrm{S}$
  & ROC-AUC & $[+0.015, +0.113]$ & $0.998$ \\
  & PR-AUC  & $[+0.006, +0.120]$ & $0.999$ \\
  & P@50    & $[-0.020, +0.120]$ & $0.850$ \\
  & P@100   & $[-0.010, +0.080]$ & $0.907$ \\
\bottomrule
\end{tabular}
\end{table}

Reading these together:
\begin{itemize}
\item The temporal-over-structural ROC-AUC gain is \emph{robust}: the
95\% CI of $(\mathrm{S}{+}\mathrm{T}) - \mathrm{S}$ is entirely above
zero, and 100\% of bootstrap resamples agree on the direction.
\item The temporal-over-structural PR-AUC and Precision@$k$ deltas on
Random Forest are \emph{directionally positive}: bootstrap means are
positive and 75--92\% of resamples are positive, but 95\% confidence
intervals touch zero. We describe these as suggestive rather than
definitive at the present test-set size.
\item The community-over-temporal PR-AUC delta on Random Forest is
also \emph{directionally positive}: bootstrap mean $+0.026$ with 94\%
of resamples favoring improvement, but the 95\% CI $[-0.003, +0.068]$
just includes zero. We therefore describe community PR-AUC gains as
suggestive, not statistically decisive.
\item The community-over-temporal ROC-AUC delta is mildly negative
on average (mean $-0.010$, 7\% of resamples positive), consistent with
saturation rather than with community features harming the model:
ROC-AUC at the $+$T level is already near $0.95$.
\item Logistic Regression shows no detectable community gain in any
metric, consistent with community signal being non-linear and not
accessible to a linear model on these features.
\item The full $\mathrm{S} \to \mathrm{S}{+}\mathrm{T}{+}\mathrm{C}$
delta on Random Forest is jointly robust on both ROC-AUC and PR-AUC
(both 95\% CIs above zero, $\sim 100\%$ of resamples positive). The
combined feature stack is meaningfully better than structural alone;
the incremental contribution of community over temporal is the part
that is only suggestive.
\item Precision@$k$ CIs are wide because Precision@$k$ is bounded by
$\min(k, n_{\mathrm{pos}})/k = 34/100 = 0.34$ for $k = 100$,
producing a hard ceiling and a skewed bootstrap distribution.
\end{itemize}

% =========================================================
\section{Error Analysis and Interpretation}
% =========================================================

We separate two kinds of evidence about community-level signal. The
first is \emph{global}, drawn from feature distributions and fitted
feature importances over the TRAIN split. The second is \emph{local},
drawn from inspection of the top-ranked test pairs from the best
model. The local analysis is descriptive and based on a very small
number of true positives; we read it as a model-behavior diagnostic,
not as confirmation or refutation of any sociological hypothesis.

\subsection{Global, feature-level evidence (TRAIN)}

Two descriptive facts at the population level are worth stating
separately:

\begin{enumerate}
\item Same-community pairs have an empirical positive rate of
$0.0140\%$ ($215 / 1{,}540{,}969$), whereas cross-community pairs are
at $0.0095\%$ ($330 / 3{,}472{,}377$). The naive ``cross-community
distrust'' reading is not supported on this dataset
(Fig.~\ref{fig:samecross}); same-community pairs are slightly
\emph{more} likely to receive a new negative rating in the next month.
\item In the fitted Logistic Regression on S$+$T$+$C,
\code{neighborhood\_overlap} is the strongest pair-level community
feature (standardized coefficient $-1.36$), in the direction that
lower shared-neighborhood predicts higher new-negative probability.
The same feature is among the top community-family signals in Random
Forest, where the dominant community-family signals are the boundary
fractions.
\end{enumerate}

Taken together, these two facts are consistent with -- but do not by
themselves prove -- a refined ``intra-community weak tie'' reading:
distrust does not predominantly appear at coarse community boundaries,
yet pair-level neighborhood overlap still matters and points in the
Granovetter direction. Establishing this reading more firmly would
require conditional analyses (e.g., overlap-conditioned positive rates
inside and across communities), which we leave to future work.

\subsection{Local diagnostic from the top-50 scored pairs (TEST, RF on
S$+$T$+$C)}

At the top-of-ranked-list operating point ($k = 50$), the best model
returns 5 true positives and 45 false positives. With $n_{TP} = 5$ the
sample is too small for statistical comparison; what follows is a
qualitative look at the kinds of pairs the model surfaces.

\begin{table}[t]
\caption{Mean feature values among the top-50 RF S$+$T$+$C scored
test pairs (illustrative; $n_{TP} = 5$, $n_{FP} = 45$).}
\label{tab:top50}
\centering
\small
\begin{tabular}{lrr}
\toprule
Feature & TPs & FPs \\
\midrule
bootstrap proba              & $0.941$ & $0.914$ \\
neg\_out\_u                  & $17.80$ & $31.53$ \\
neg\_in\_v                   & $25.60$ & $14.56$ \\
decayed\_neg\_received\_v    & $1.34$  & $2.04$  \\
days\_since\_last\_activity\_u & $3.59$ & $2.53$ \\
same\_community              & $0.00$  & $0.16$  \\
neighborhood\_overlap        & $0.080$ & $0.055$ \\
boundary\_frac\_u            & $0.295$ & $0.394$ \\
uv\_interaction\_count\_past & $0.00$  & $0.00$  \\
\bottomrule
\end{tabular}
\end{table}

Three qualitative patterns are visible
(Table~\ref{tab:top50}). First, the model concentrates risk on a
small set of recurring ``issuer'' nodes (high \code{neg\_out\_u}) and
recurring ``victim'' nodes (high \code{neg\_in\_v}); this matches the
features. Second, what distinguishes the small true-positive subset
from the false positives at the very top is mostly on the
\emph{victim} side (mean \code{neg\_in\_v} 25.6 vs 14.6 for FPs), not
the issuer side. False positives at the top tend to be highly active
negative-givers whose specific $(u, v)$ target did not receive a new
negative in the test window. Third, on the community-aware features
the local direction in this slice does \emph{not} mirror the global
TRAIN pattern: in the top-50, all 5 TPs are cross-community, and the
mean \code{neighborhood\_overlap} is higher (not lower) for TPs than
for FPs.

We deliberately do not interpret the top-50 slice as confirming or
refuting the global pattern. With 5 true positives, the local
direction is sensitive to which specific candidate pairs the model
happened to rank highly, and the local statistic does not have the
sample support to override the population-level evidence from the
previous subsection. Supplementary top-200 and top-500 slices
(available on disk as \code{results/stage7\_error\_top200.csv} and
\code{results/stage7\_error\_top500.csv}) further illustrate that
the model is highly \emph{front-loaded}: only 7 of the top 200 and
8 of the top 500 scored pairs are true positives. Most of the model's
correct mass is concentrated in a small number of high-confidence
recurring patterns.

\subsection{Summary}

The descriptive evidence on TRAIN is consistent with a refined
weak-tie reading in which lower pair-level neighborhood overlap is
associated with higher distrust risk, while cross-community status by
itself is not. The local top-50 diagnostic shows that the model's
highest-confidence predictions are concentrated on recurring victims
and that its top false positives are dominated by highly active
negative-givers; the community-feature direction in this small slice
does not align with the global feature-importance direction, which is
exactly the kind of conditional-vs-marginal divergence one expects
when the sample is this small.

% =========================================================
\section{Limitations}\label{sec:lim}
% =========================================================

We list the limitations explicitly so that follow-up work can
prioritize.

\begin{itemize}
\item \textbf{Few positives in test.} The locked test split has 34
positive candidates. This is the structural source of wide PR-AUC and
Precision@$k$ confidence intervals; with this many positives the
paired stratified bootstrap cannot distinguish small community-feature
effects from sampling noise.
\item \textbf{Single network, no external replication.} We do not
report results on Bitcoin-Alpha or any other signed network. A
natural extension is zero-shot transfer to Bitcoin-Alpha; we leave
this as future work.
\item \textbf{Single time window for test.} TEST is exactly two
months; rolling-origin evaluation across many windows would yield a
larger effective positive count and tighter intervals.
\item \textbf{Default hyperparameters.} We do not sweep model or
feature thresholds (\code{MIN\_ACTIVITY}, \code{RECENT\_DAYS},
\code{MAX\_CANDIDATES}, \code{K\_HOP}). The intent is interpretable
analysis at locked defaults.
\item \textbf{Last-rating-wins snapshot collapse.} The snapshot graph
uses the latest rating per pair; this represents current trust
judgment but discards multi-event nuance. We restore the full event
history only for temporal and trust-summary features.
\item \textbf{Subsampled training.} Negatives in TRAIN are subsampled
$100{:}1$ for tractability. Validation and test remain full and
untouched. A small calibration loss is possible at this ratio but does
not affect ranking metrics in our setup.
\item \textbf{Percentile CIs.} We report percentile bootstrap CIs,
which can under-cover when the bootstrap distribution is skewed -- in
particular for Precision@$k$ near the hard ceiling. BCa or
studentized intervals would be a natural strengthening.
\item \textbf{Iterated fairness/goodness deferred.} Iterated
fairness and goodness in the sense of~\cite{kumar2016} is not run; our
four simple trust-summary features are a deliberately lighter
analogue.
\item \textbf{\code{prior\_interaction} is partly degenerate.} By
construction of the candidate set, this flag is $1$ only when $u$
previously rated $v$ positively. We retain it for transparency but
note that its fitted coefficient should be read with care.
\item \textbf{Random Forest \texttt{predict\_proba} non-determinism.}
With \code{n\_jobs = -1}, RF probabilities can differ at the
float-tie boundary between two independent fits. This produced a
$\sim 0.01$ PR-AUC drift between the original Stage~7 metrics and the
canonical metrics in Table~\ref{tab:test_metrics}; the drift is well
inside the bootstrap CI and does not change any qualitative
conclusion.
\item \textbf{Compute timing not reported.} We do not include a
computational-cost table; persisted runtime logs were not available
and we did not rerun the pipeline to recover them. \todo{add timing
table if rerun is acceptable before submission.}
\item \textbf{\emph{frac\_boot}$(>0)$ is not a p-value.} We use it
only to summarize the direction of the bootstrap distribution; it
carries no formal frequentist coverage guarantee.
\end{itemize}

% =========================================================
\section{Conclusion}
% =========================================================

We presented an interpretable, leakage-controlled pipeline for early
warning of new negative trust-edge emergence in the Bitcoin-OTC
temporal signed network, together with a three-way feature ablation
under Logistic Regression and Random Forest. On a locked test split
of 34 positives over 148{,}639 candidate pairs, structural features
alone already produce a non-trivial ranker (Random Forest ROC-AUC
$0.88$, PR-AUC $0.022$). Adding temporal features improves ranking
robustly, particularly for Random Forest; the
$(\mathrm{S}{+}\mathrm{T}) - \mathrm{S}$ ROC-AUC delta sits entirely
above zero on every one of 1{,}000 paired bootstrap resamples.

Adding community-aware features on top of structural and temporal
features tells a more nuanced story. Community features do not
improve ROC-AUC over the temporal baseline; on Random Forest the
ROC-AUC delta is mildly negative on average, consistent with a
saturated ceiling. They do shift Random Forest PR-AUC and
Precision@50 in the positive direction (bootstrap mean
$+0.026$ PR-AUC, 94\% of resamples positive), but the 95\% CIs still
include zero, so we describe community gains as suggestive rather
than statistically decisive at this test-set size. Logistic
Regression shows no community gain in any metric, consistent with the
community signal being non-linear.

At the descriptive level, the data are not consistent with a coarse
``distrust between communities'' reading: same-community pairs in
TRAIN have a higher empirical positive rate than cross-community
pairs. They are consistent with a finer pair-level pattern in which
lower neighborhood overlap is associated with higher new-negative
probability. We frame this as a refined, conservative reading rather
than as a confirmed weak-tie result; a stronger claim would require
overlap-conditioned analyses inside and across communities that we
leave to future work.

The largest natural extensions are external replication on
Bitcoin-Alpha, rolling-origin evaluation across more test windows so
that the community-effect intervals can draw on more positives, and
BCa or studentized bootstrap CIs to address skew at the
Precision@$k$ ceiling.

% =========================================================
% REFERENCES
% =========================================================
\begin{thebibliography}{99}

\bibitem{kumar2016}
S.~Kumar, F.~Spezzano, V.~S. Subrahmanian, and C.~Faloutsos,
``Edge weight prediction in weighted signed networks,''
in \emph{Proc.\ IEEE 16th Int.\ Conf.\ Data Mining (ICDM)},
Barcelona, Spain, Dec.\ 2016, pp.\ 221--230.

\bibitem{bertazzi2018}
I.~Bertazzi, S.~Huet, G.~Deffuant, and F.~Gargiulo,
``The anatomy of a web of trust: The Bitcoin-OTC market,''
in \emph{Social Informatics (SocInfo 2018)},
Lecture Notes in Computer Science, vol.\ 11185, Springer, 2018.
DOI: 10.1007/978-3-030-01129-1\_14.
\todo{verify page numbers within LNCS 11185 before camera-ready.}

\bibitem{choudhury2024}
N.~Choudhury,
``Community-aware evolution similarity for link prediction in dynamic
social networks,''
\emph{Mathematics (MDPI)}, vol.\ 12, no.\ 2, art.\ 285, 2024.
DOI: 10.3390/math12020285.

\bibitem{zhang2024}
Z.~Zhang, P.~Zhao, X.~Li, J.~Liu, X.~Zhang, J.~Huang, and X.~Zhu,
``Signed graph representation learning: A survey,''
\emph{arXiv preprint arXiv:2402.15980}, Feb.\ 2024.

\bibitem{granovetter1973}
M.~S. Granovetter,
``The strength of weak ties,''
\emph{Amer.\ J.\ Sociology}, vol.\ 78, no.\ 6, pp.\ 1360--1380, 1973.
DOI: 10.1086/225469.

\bibitem{blondel2008}
V.~D. Blondel, J.-L. Guillaume, R.~Lambiotte, and E.~Lefebvre,
``Fast unfolding of communities in large networks,''
\emph{J.\ Stat.\ Mech.: Theory Exp.}, vol.\ 2008, no.\ 10, art.\ P10008,
Oct.\ 2008. DOI: 10.1088/1742-5468/2008/10/P10008.

\bibitem{adamic2003}
L.~A. Adamic and E.~Adar,
``Friends and neighbors on the web,''
\emph{Social Networks}, vol.\ 25, no.\ 3, pp.\ 211--230, 2003.
\todo{confirm DOI 10.1016/S0378-8733(03)00009-1 manually before
camera-ready.}

\bibitem{snap}
J.~Leskovec and A.~Krevl,
``SNAP datasets: Stanford large network dataset collection,''
\url{http://snap.stanford.edu/data}, Jun.\ 2014.
Bitcoin-OTC subset:
\url{https://snap.stanford.edu/data/soc-sign-bitcoinotc.html}.

\end{thebibliography}

\end{document}
